% =============================================================================
% APPENDIX PT: DOMAIN-PLATFORM: Time-Intensive Goods & Platform Economics
% =============================================================================
% Module: BCM2_17_PLATFORM
% Version: 1.0 (January 2026)
% Status: Final
%
% This appendix applies Becker's Time Allocation Theory to digital platform
% economics, analyzing consumer demand and market competition for time-intensive
% goods. Based on Goodman et al. (2026) NBER WP 34743.
%
% =============================================================================

% =============================================================================
% CROSS-REFERENCE STRUCTURE
% =============================================================================
%
% CHAPTER LINKAGE:
% - Primary Chapter: Chapter 9 (Context and the Ψ-Framework)
% - Secondary Chapters: Chapter 17 (Interventions), Chapter 20 (Antitrust)
%
% DEPENDENCIES (what this appendix REQUIRES):
% - Appendix V (CORE-WHEN): Temporal context and Ψ_T dimension
% - Appendix AAA (CORE-WHO): Consumer segments and levels
% - Appendix C (CORE-WHAT): Utility dimensions (FEPSDE)
% - Appendix BBB (CORE-WHERE): Parameter values
% - Appendix MS (REF-MODELSPACE): Theory catalog CAT-17
%
% DEPENDENTS (what REQUIRES this appendix):
% - Appendix DER (FORMAL-LIMITS): Proofs for §7.1, §7.2
% - Antitrust applications (future)
% - Digital platform interventions
%
% =============================================================================

\chapter{Time-Intensive Goods \& Platform Economics}
\label{app:platform-economics}

% -----------------------------------------------------------------------------
% HEADER BLOCK (Appendix Metadata)
% -----------------------------------------------------------------------------

\begin{tcolorbox}[colback=gray!5!white, colframe=gray!75!black,
    title=Appendix Metadata]
\begin{tabular}{@{}ll@{}}
\textbf{Code:} & PT \\
\textbf{Category:} & DOMAIN (Domain Applications) \\
\textbf{Name:} & PLATFORM \\
\textbf{Title:} & Time-Intensive Goods \& Platform Economics \\
\textbf{Version:} & 1.0 (January 2026) \\
\textbf{Status:} & ACTIVE \\
\textbf{Dependencies:} & V, AAA, C, BBB, MS \\
\textbf{Primary Chapter:} & Chapter 9 (Context and the $\Psi$-Framework) \\
\textbf{Secondary Chapters:} & Chapter 17, Chapter 20 \\
\end{tabular}
\end{tcolorbox}

% -----------------------------------------------------------------------------
% CHAPTER LINKAGE BOX
% -----------------------------------------------------------------------------

\begin{tcolorbox}[colback=purple!5!white, colframe=purple!75!black,
    title=Chapter Linkage]

\textbf{Primary Chapter:} Chapter 9: Context and the $\Psi$-Framework
\begin{itemize}[nosep]
\item Section 9.3: Temporal Context ($\Psi_T$) $\leftrightarrow$ This Appendix Section \ref{sec:pt-time-constraint}
\item Section 9.5: Context Interactions $\leftrightarrow$ This Appendix Section \ref{sec:pt-platform-application}
\end{itemize}

\textbf{Secondary Chapters:}
\begin{itemize}[nosep]
\item Chapter 17: Intervention Design --- attention nudges, time salience
\item Chapter 20: Policy Applications --- antitrust, market definition
\end{itemize}

\textbf{Reading Guide:}
\begin{itemize}[nosep]
\item For conceptual overview: Read Chapter 9 first (Context Framework)
\item For formal details: This appendix provides complete Beckerian specification
\item For proofs: See Appendix DER §7.1, §7.2
\end{itemize}

\end{tcolorbox}

% -----------------------------------------------------------------------------
% ABSTRACT
% -----------------------------------------------------------------------------

\begin{abstract}
This appendix integrates Becker's (1965) Time Allocation Theory into the EBF framework,
with particular application to digital platform economics. We formalize how time constraints
affect consumer demand when monetary prices approach zero, deriving predictions for
substitution patterns (diversion ratios) based on time cost shares. The framework explains
why Facebook and Instagram are close substitutes despite being ``free''---their similar
time shares predict high diversion ratios. Policy implications include revised antitrust
market definitions and attention-based intervention design.

\textit{Theory Catalog Reference: CAT-17 (MS-TA-001, MS-TA-002, MS-TA-003)}
\end{abstract}

% -----------------------------------------------------------------------------
% QUICK REFERENCE BOX
% -----------------------------------------------------------------------------

\begin{tcolorbox}[colback=blue!5!white, colframe=blue!75!black,
    title=Quick Reference: Time Allocation \& Platform Economics]

\textbf{Core Equations:}
\begin{align}
\text{Full Price:} \quad & \pi_i = p_i + W \cdot t_i \label{eq:pt-full-price} \\
\text{Time Share:} \quad & s_T^i = \frac{W \cdot t_i}{\pi_i} = \frac{W \cdot t_i}{p_i + W \cdot t_i} \label{eq:pt-time-share} \\
\text{Diversion Ratio:} \quad & D_{ij} = \frac{\partial q_j / \partial p_i}{-\partial q_i / \partial p_i} \label{eq:pt-diversion}
\end{align}

\textbf{Key Parameters (from PAR-TA-001 to PAR-TA-005):}
\begin{center}
\begin{tabular}{llll}
\toprule
Symbol & Name & Typical Value & Source \\
\midrule
$W$ & Shadow wage & CHF 42.50/h (CH) & BFS 2024 \\
$s_T$ & Time share (social media) & $\approx 0.99$ & Goodman et al. 2026 \\
$D_{ij}$ & FB-IG Diversion & $\approx 0.65$ & Goodman et al. 2026 \\
$\varepsilon_{ad}$ & Ad elasticity & $\approx -0.15$ & Goodman et al. 2026 \\
$A_{max}$ & Daily attention budget & 720 min & Simon 1971 \\
\bottomrule
\end{tabular}
\end{center}

\textbf{Beckerian Prediction:}
\[
D_{ij} \uparrow \quad \text{when} \quad |s_T^i - s_T^j| \downarrow
\]
Products with similar time shares are closer substitutes.

\end{tcolorbox}

% -----------------------------------------------------------------------------
% APPENDIX SCOPE BOX
% -----------------------------------------------------------------------------

\begin{tcolorbox}[colback=orange!5!white, colframe=orange!75!black,
    title=Appendix Scope: Time-Intensive Goods \& Platform Economics]

\textbf{Ziel:} Provide formal EBF integration of time allocation theory for platform economics and antitrust analysis.

\textbf{In-Scope:}
\begin{itemize}[nosep]
\item Becker (1965) time allocation model and EBF restrictions
\item Time share formula and diversion ratio predictions
\item Digital platform application (Facebook/Instagram case)
\item Antitrust implications (market definition, merger analysis)
\item Attention economy connection (Simon 1971)
\end{itemize}

\textbf{Out-of-Scope:}
\begin{itemize}[nosep]
\item Detailed welfare analysis (see Appendix WEL)
\item Network effects modeling (see Appendix NET)
\item Platform pricing strategies (advertising side)
\item Regulatory implementation details
\end{itemize}

\textbf{Deliverables:}
\begin{itemize}[nosep]
\item Formal axioms (TA-1 to TA-5)
\item Parameter estimates (PAR-TA-001 to PAR-TA-005)
\item Worked example (Facebook-Instagram substitution)
\item Policy implications checklist
\end{itemize}

\end{tcolorbox}

% =============================================================================
% PART 2: CORE CONTENT
% =============================================================================

\section{The Fundamental Question}
\label{sec:pt-fundamental}

\begin{tcolorbox}[colback=yellow!10!white, colframe=orange!75!black,
    title=Central Question]
\textbf{How do consumers allocate scarce time across competing activities, and what does this imply for market competition when monetary prices are zero or negligible?}
\end{tcolorbox}

Traditional demand analysis assumes consumers face budget constraints measured in money. But for ``free'' digital platforms like Facebook, Instagram, TikTok, and YouTube, the binding constraint is \textit{time}, not money. This fundamentally changes substitution patterns:

\begin{itemize}
\item \textbf{Traditional goods:} Competition on price $\Rightarrow$ SSNIP test (5\% price increase)
\item \textbf{Time-intensive goods:} Competition on attention $\Rightarrow$ Time-based market definition
\end{itemize}

% -----------------------------------------------------------------------------
\section{Becker's Time Allocation Theory}
\label{sec:pt-time-constraint}

\subsection{The Household Production Function}

Following \citet{PAP-becker1965theory}, consumers do not derive utility directly from market goods $X$ but from \textit{commodities} $Z$ produced using goods and time:

\begin{axiom}[TA-1: Household Production]
\label{ax:ta-household}
Utility is derived from commodities $Z_i$ produced via household production functions:
\[
Z_i = f_i(X_i, T_i)
\]
where $X_i$ is market goods input and $T_i$ is time input.
\end{axiom}

\subsection{The Time Constraint}

\begin{axiom}[TA-2: Time Constraint]
\label{ax:ta-time}
Total time $T$ is fixed and allocated across activities:
\[
T = T_{work} + \sum_i T_i + T_{leisure}
\]
This constraint is always binding---unlike money, time cannot be borrowed or saved.
\end{axiom}

\subsection{The Full Price}

\begin{axiom}[TA-3: Full Price]
\label{ax:ta-full-price}
The full price $\pi_i$ of commodity $Z_i$ includes both monetary cost and time cost:
\[
\pi_i = p_i + W \cdot t_i
\]
where:
\begin{itemize}[nosep]
\item $p_i$ = monetary price of goods input
\item $W$ = shadow price of time (opportunity cost, $\approx$ wage rate)
\item $t_i$ = time required per unit of $Z_i$
\end{itemize}
\end{axiom}

\subsection{The Time Share}

\begin{definition}[Time Cost Share]
The time share $s_T^i$ measures the fraction of full price attributable to time:
\[
s_T^i = \frac{W \cdot t_i}{\pi_i} = \frac{W \cdot t_i}{p_i + W \cdot t_i}
\]
\end{definition}

\begin{proposition}[Limiting Cases]
\label{prop:pt-limits}
\begin{enumerate}[nosep]
\item If $p_i = 0$ (free good): $s_T^i \to 1$ (pure time cost)
\item If $t_i = 0$ (instant consumption): $s_T^i \to 0$ (pure monetary cost)
\item For social media platforms: $p_i \approx 0 \Rightarrow s_T \approx 1$
\end{enumerate}
\end{proposition}

% -----------------------------------------------------------------------------
\section{Diversion Ratios and Substitution}
\label{sec:pt-diversion}

\subsection{The Beckerian Prediction}

\begin{axiom}[TA-4: Time-Based Substitution]
\label{ax:ta-substitution}
Products with similar time shares are closer substitutes:
\[
\frac{\partial D_{ij}}{\partial |s_T^i - s_T^j|} < 0
\]
where $D_{ij}$ is the diversion ratio from product $i$ to product $j$.
\end{axiom}

\begin{definition}[Diversion Ratio]
The diversion ratio $D_{ij}$ measures the proportion of lost demand from product $i$ that diverts to product $j$:
\[
D_{ij} = \frac{\partial q_j / \partial p_i}{-\partial q_i / \partial p_i}
\]
where $\sum_j D_{ij} \leq 1$ (some demand may exit the market).
\end{definition}

\subsection{Empirical Findings}

\citet{PAP-goodman2026timeintensive} conducted two field experiments on Facebook and Instagram users, finding:

\begin{tcolorbox}[colback=green!5!white, colframe=green!75!black,
    title=Key Empirical Results (Goodman et al. 2026)]
\begin{enumerate}
\item \textbf{High FB-IG Diversion:} $D_{FB \to IG} \approx 0.65$ [CI: 0.55--0.75]
\item \textbf{Low Ad Elasticity:} $\varepsilon_{ad} \approx -0.15$ [CI: -0.30 to -0.05]
\item \textbf{Time Shares Near Unity:} $s_T^{FB} \approx s_T^{IG} \approx 0.99$
\item \textbf{Beckerian Prediction Confirmed:} Similar $s_T$ predicts high $D_{ij}$
\end{enumerate}
\end{tcolorbox}

% -----------------------------------------------------------------------------
\section{The Attention Economy}
\label{sec:pt-attention}

\begin{axiom}[TA-5: Attention Scarcity]
\label{ax:ta-attention}
Following \citet{PAP-simon1971designing}, attention is bounded:
\[
A = \sum_i a_i \leq A_{max}
\]
where $A_{max}$ is the daily attention budget ($\approx$ 720 minutes of discretionary time).

\textbf{Simon's Insight:} ``A wealth of information creates a poverty of attention.''
\end{axiom}

This connects time allocation theory to the broader attention economy:
\begin{itemize}
\item Platforms compete for attention, not money
\item User ``payment'' is in time and attention
\item Ad load is the implicit price users pay
\end{itemize}

% -----------------------------------------------------------------------------
\section{EBF Integration}
\label{sec:pt-ebf-integration}

\subsection{10C Mapping}

\begin{center}
\begin{tabular}{lll}
\toprule
10C Dimension & Relevance & Time Allocation Component \\
\midrule
\textbf{WHEN} ($\Psi_T$) & Primary & Time constraint $T$, shadow wage $W$ \\
\textbf{WHAT} (C) & Primary & Utility from commodities $Z_i$ \\
\textbf{WHO} (L) & Secondary & Consumer segments by $W$ \\
\textbf{WHERE} ($\Theta$) & Secondary & Parameters $s_T$, $D_{ij}$, $\varepsilon_{ad}$ \\
\textbf{AWARE} (A) & Extended & Awareness of time opportunity cost \\
\bottomrule
\end{tabular}
\end{center}

\subsection{Theory Catalog Connection}

This appendix implements theories from CAT-17 (Household Economics \& Time Allocation):

\begin{itemize}
\item \textbf{MS-TA-001:} Time Allocation Theory (Becker 1965) --- foundational
\item \textbf{MS-TA-002:} Time-Intensive Goods \& Market Competition (Goodman et al. 2026) --- digital platform extension
\item \textbf{MS-TA-003:} Attention Economy (Simon 1971) --- attention scarcity
\end{itemize}

% =============================================================================
% PART 3: APPLICATION
% =============================================================================

\section{Worked Example: Facebook-Instagram Substitution}
\label{sec:pt-platform-application}

\begin{tcolorbox}[colback=blue!5!white, colframe=blue!75!black,
    title=Case Study: CAS-853]

\textbf{Context:} The FTC has proposed a de-merger of Facebook and Instagram, arguing they are close substitutes. How can we quantify this using time allocation theory?

\textbf{Given:}
\begin{itemize}[nosep]
\item Both platforms are free: $p_{FB} = p_{IG} \approx 0$
\item Average user time: $t_{FB} \approx 35$ min/day, $t_{IG} \approx 30$ min/day
\item Swiss median wage: $W = $ CHF 42.50/hour
\end{itemize}

\textbf{Step 1: Calculate Time Shares}
\[
s_T^{FB} = \frac{W \cdot t_{FB}}{0 + W \cdot t_{FB}} = 1.0
\]
\[
s_T^{IG} = \frac{W \cdot t_{IG}}{0 + W \cdot t_{IG}} = 1.0
\]

\textbf{Step 2: Apply Beckerian Prediction}

Since $|s_T^{FB} - s_T^{IG}| \approx 0$, the theory predicts high diversion.

\textbf{Step 3: Empirical Validation}

Goodman et al. (2026) found $D_{FB \to IG} \approx 0.65$, confirming the prediction.

\textbf{Conclusion:} Facebook and Instagram are close substitutes in the time-intensive goods sense. A de-merger could increase competition in the attention market.

\end{tcolorbox}

% -----------------------------------------------------------------------------
\section{Policy Implications}
\label{sec:pt-policy}

\subsection{Antitrust Market Definition}

\begin{tcolorbox}[colback=red!5!white, colframe=red!75!black,
    title=Implication 1: SSNIP Test Failure]
Traditional SSNIP tests (5\% price increase) fail when $p \approx 0$:
\begin{itemize}[nosep]
\item 5\% of zero is still zero
\item Need \textbf{time-based} market definition
\item Relevant question: ``Would a 5\% increase in ad load cause switching?''
\end{itemize}
\end{tcolorbox}

\subsection{Platform Market Power}

\begin{tcolorbox}[colback=red!5!white, colframe=red!75!black,
    title=Implication 2: Ad Load Pricing Power]
Low ad elasticity ($\varepsilon_{ad} \approx -0.15$) indicates:
\begin{itemize}[nosep]
\item Users tolerate high ad loads
\item Platforms have significant market power on user side
\item Regulation may be justified even without monetary pricing
\end{itemize}
\end{tcolorbox}

\subsection{Intervention Design}

\begin{tcolorbox}[colback=red!5!white, colframe=red!75!black,
    title=Implication 3: Attention Nudges]
EBF interventions should target time awareness:
\begin{itemize}[nosep]
\item Make time opportunity cost salient (``You've spent 2 hours---worth CHF 85'')
\item Default screen time limits
\item Activity reports showing time allocation
\end{itemize}
\end{tcolorbox}

% =============================================================================
% PART 4: BACK MATTER
% =============================================================================

\section{Summary}
\label{sec:pt-summary}

\begin{tcolorbox}[colback=gray!10!white, colframe=gray!75!black,
    title=Key Takeaways]
\begin{enumerate}
\item \textbf{Time as Currency:} For free platforms, time is the binding constraint, not money.
\item \textbf{Full Price:} $\pi = p + W \cdot t$ includes both monetary and time costs.
\item \textbf{Time Share:} $s_T = (W \cdot t)/\pi$ determines substitution patterns.
\item \textbf{Beckerian Prediction:} Similar $s_T \Rightarrow$ high diversion ratios.
\item \textbf{Platform Implications:} FB-IG diversion $\approx 0.65$ supports de-merger argument.
\item \textbf{Policy:} Traditional antitrust tools need time-based adaptation.
\end{enumerate}
\end{tcolorbox}

% -----------------------------------------------------------------------------
\section{Glossary of Symbols}
\label{sec:pt-glossary}

See \textbf{Appendix G (Glossary)} for comprehensive definitions. Key symbols in this appendix:

\begin{center}
\begin{tabular}{lll}
\toprule
Symbol & Name & Definition \\
\midrule
$\pi_i$ & Full price & $p_i + W \cdot t_i$ \\
$s_T^i$ & Time share & $(W \cdot t_i) / \pi_i$ \\
$W$ & Shadow wage & Opportunity cost of time \\
$D_{ij}$ & Diversion ratio & $(\partial q_j/\partial p_i) / (-\partial q_i/\partial p_i)$ \\
$A_{max}$ & Attention budget & Daily discretionary time \\
$\varepsilon_{ad}$ & Ad elasticity & $\partial \ln q / \partial \ln (\text{ad load})$ \\
$\Psi_T$ & Temporal context & EBF context dimension for time \\
\bottomrule
\end{tabular}
\end{center}

% -----------------------------------------------------------------------------
\section{Critical Foundations}
\label{sec:pt-foundations}
% -----------------------------------------------------------------------------

\begin{enumerate}
\item \textbf{Becker (1965):} Time allocation theory provides the foundational insight that consumption requires time as well as money. The full price $\pi = p + Wt$ integrates both.

\item \textbf{Simon (1971):} Attention is the scarce resource in an information-rich world. This insight grounds the attention economy framework.

\item \textbf{Goodman et al. (2026):} Empirical validation of Beckerian predictions for digital platforms. Time share determines substitution patterns.

\item \textbf{EBF Integration:} The temporal context dimension $\Psi_T$ (Appendix V) captures time scarcity effects. Platform economics is a domain application.

\item \textbf{Antitrust Economics:} Market definition based on time costs rather than monetary prices. Products with similar $s_T$ are close substitutes.
\end{enumerate}

% -----------------------------------------------------------------------------
\section{Open Issues}
\label{sec:pt-open}

\begin{enumerate}
\item \textbf{Heterogeneous Shadow Wages:} How do substitution patterns vary by income/education?
\item \textbf{Network Effects:} How do social network effects interact with time allocation?
\item \textbf{Multi-Homing:} Many users use both FB and IG---how to model joint consumption?
\item \textbf{Welfare Measurement:} How to measure consumer surplus when $p=0$?
\item \textbf{Dynamic Effects:} How does time allocation change over the life cycle?
\end{enumerate}

% -----------------------------------------------------------------------------
\section{References}
\label{sec:pt-references}

This appendix draws on the \textbf{Master Bibliography} (bcm\_master.bib). Key references:

\begin{itemize}[nosep]
\item \citet{PAP-becker1965theory}: Foundational time allocation theory
\item \citet{PAP-goodman2026timeintensive}: Digital platform application (NBER WP 34743)
\item \citet{PAP-simon1971designing}: Attention scarcity insight
\end{itemize}

\nocite{bcm_master}

% -----------------------------------------------------------------------------
% CROSS-REFERENCE FOOTER
% -----------------------------------------------------------------------------

\begin{tcolorbox}[colback=white, colframe=black,
    title=Cross-References]

\textbf{Related Appendices:}
\begin{itemize}[nosep]
\item Appendix V (CORE-WHEN): Temporal context $\Psi_T$
\item Appendix AAA (CORE-WHO): Consumer segmentation by shadow wage
\item Appendix BBB (CORE-WHERE): Parameter values PAR-TA-001 to PAR-TA-005
\item Appendix MS (REF-MODELSPACE): Theory catalog CAT-17
\item Appendix DER (FORMAL-LIMITS): Formal proofs §7.1, §7.2
\end{itemize}

\textbf{Related Models:}
\begin{itemize}[nosep]
\item MOD-TA-001: Time-Intensive Goods Substitution Model
\end{itemize}

\textbf{Related Cases:}
\begin{itemize}[nosep]
\item CAS-853: Facebook-Instagram Zeit-Substitution
\end{itemize}

\end{tcolorbox}

% =============================================================================
% END OF APPENDIX PT
% =============================================================================
